\documentclass[11pt,twoside,openany,full]{on1x-book}
\OnxBookTitle{On1x: A Language and Its Toolchain}
\OnxBookAuthor{Your Name}
\OnxBookDate{Draft edition, November 2026}
\OnxPDFMetadata{On1x: A Language and Its Toolchain}{Your Name}{Technical programming language book}
\OnxGlossaryEntry{hir}{HIR}{High-level intermediate representation}
\OnxAcronym{mir}{MIR}{mid-level intermediate representation}
\begin{document}
\frontmatter
\maketitle
\OnxCopyrightPage{2026}{Your Name}{Replace with your own licensing terms and edition statement.}
\OnxContents
\mainmatter
\part{The Language}
\chapter{Introducing On1x}\label{ch:intro}
\OnxChapterLead{A statically compiled language, its standard library, and its C11 toolchain.}
\begin{OnxObjectives}
\item Distinguish source syntax from the implementation status of a feature.
\item Read code snippets, formal judgments, and compiler diagrams.
\item Identify evidence and capability limitations.
\end{OnxObjectives}
\section{A sample of the notation}
The language is called \OnxLanguage{}, and its kernel ecosystem is \OnxMuLIB{}.
Inline examples such as \OnxInline|fn| and \OnxInline|let| are typeset as code.
The source below is \emph{illustrative syntax}, not an execution claim.
\begin{OnxCodeBlock}[caption={Illustrative function syntax},label={lst:first}]
fn twice(x: i64) -> i64 {
    return x + x;
}
\end{OnxCodeBlock}
\begin{OnxNote}
Use \OnxPartial{} or \OnxAbsent{} for planned features rather than claiming that a code sample is executable.
\end{OnxNote}
\begin{OnxWarning}
Feature descriptions and kernel headers are not evidence that an On1x binding is admitted or safe.
\end{OnxWarning}
\section{An explanatory example}
\begin{OnxExample}[A three-stage compilation outline]\label{ex:stages}
An implementation may parse source, check a high-level intermediate form, and lower toward a target.
\end{OnxExample}
\begin{OnxDiagram}
\OnxDiagramNode{src}{Source}
\OnxDiagramNode[right=of src]{hir}{HIR}
\OnxDiagramNode[right=of hir]{c}{C11}
\OnxDiagramArrow{src}{hir}
\OnxDiagramArrow{hir}{c}
\end{OnxDiagram}
\begin{OnxExercise}[Read the compiler structure][Easy]\label{ex:compiler}
Explain why a checked intermediate representation is useful.
\end{OnxExercise}
\begin{OnxSolution}
It provides a point for verification and optimization before target lowering.
\end{OnxSolution}
\section{Formal notation}
Typing judgments can be written as
\[
  \OnxTyping{\Gamma}{e}{\tau}
\]
The following example is a schema, not a normative On1x inference rule:
\begin{OnxRules}
  \OnxRule[Demo]{\Gamma \vdash e : \tau}{\Gamma \vdash e : \tau}
\end{OnxRules}
\begin{OnxDefinition}[Demo type]\label{def:type}
A type classifies a set of values in a particular static semantics.
\end{OnxDefinition}
\begin{OnxTheorem}[Reflexivity of equality]
For any integer $n$, $n=n$.
\end{OnxTheorem}
\begin{proof}
Equality is reflexive.
\end{proof}
\section{Feature status is version-specific}
\begin{OnxFeatureMatrix}
\OnxFeatureRow{Cooperative channels}{\OnxExecutable}{Scoped profile}
\OnxFeatureRow{Async/await}{\OnxAbsent}{Not implemented in checked source snapshot}
\end{OnxFeatureMatrix}
See \OnxRef{ch:intro} and \OnxRef{ex:stages} for cross-reference formatting.
\begin{OnxSummary}
\item Write examples conservatively and describe the source revision.
\item Keep reference notation consistent across chapters.
\end{OnxSummary}
\chapter{Next Steps}
\begin{OnxShell}
$ pdflatex main.tex
$ bibtex main
$ pdflatex main.tex
\end{OnxShell}
\begin{OnxProject}{Build a tiny language reference}
\OnxDeliverable{A chapter explaining one language construct with syntax and semantics.}
\OnxAcceptance{Every execution claim includes a profile and a test reference.}
\end{OnxProject}
\backmatter
\OnxAppendixChapter{Notation Reference}
\begin{OnxNotationTable}
\OnxNotationRow{$\Gamma$}{A typing context}
\OnxNotationRow{$\tau$}{A type}
\end{OnxNotationTable}
\OnxPrintIndex
\end{document}
